%! Characteristics of Quadratic Functions — Unit 2, Lesson 2.0 — Homework (Answer Key)
% Mirrors homework/main.tex EXACTLY: -key in place of -boxes, every \blank{} becomes an \ans{},
% every \writelines{n} is answered with exactly n \ansline{}s, and the one work block (item 5a)
% is authored byte-identically, so the two files come out the same number of pages. No teacher
% notes here — they live in the plan. Trimmed in lockstep with the blank on 2026-09-17.
\documentclass[12pt]{article}
\usepackage{algebra2-article}
\usepackage{algebra2-key}   % pulls in -boxes; do NOT also load -boxes
% Coordinate grid: {xmin}{xmax}{ymin}{ymax}
\newcommand{\lgrid}[4]{%
  \draw[step=1, linegray!40, very thin] (#1,#3) grid (#2,#4);
  \draw[->, semithick, charcoal] (#1,0)--(#2,0) node[below right, font=\tiny]{$x$};
  \draw[->, semithick, charcoal] (0,#3)--(0,#4) node[left, font=\tiny]{$y$};}
\newcommand{\fdot}[3]{\fill[#1] (#2,#3) circle (2.6pt);}

\begin{document}

\pageheader{Unit 2, Lesson 2.0}{Homework --- Answer Key}
% NAMESTRIP: no \namedateperiod here — mirrors the blank.

\boxguard[8]
\begin{remindbox}
\small
\textbf{This is your graded homework.} Your packet with completed homework is due the first class
after two study halls. I will announce the due date in class and on TurtleNet. Remember: the
\textbf{maximum} or \textbf{minimum value} is an \emph{output} --- a $y$-value, with units --- and
\emph{where} it happens is an \emph{input}, the $x$-coordinate of the vertex.
\end{remindbox}

\vspace{0.08in}

\boxguard[14]
\begin{notesbox}{Practice}
\small
\begin{enumerate}[leftmargin=1.9em, itemsep=7pt]
  \item \textbf{Read the whole graph.} Below is $f(x)=x^2+2x-3$.
    \begin{minipage}[c]{0.68\linewidth}
    Opens \ans{up} \; Vertex \ans{$(-1,-4)$} \; Axis \ans{$x=-1$} \\[2pt]
    Zeros \ans{$-3$, $1$} \; $y$-intercept \ans{$(0,-3)$} \\[2pt]
    Increasing on \ans{$[-1,\infty)$} \; Decreasing on \ans{$(-\infty,-1]$} \\[2pt]
    Domain \ans{all real numbers} \; Range \ans{$y\ge-4$} \\[2pt]
    The vertex is a \ans{minimum} point; that \textbf{value} is \ans{$-4$}, at $x=$
    \ans{$-1$}.
    \end{minipage}\hfill
    \begin{minipage}[c]{0.28\linewidth}\centering
    \begin{tikzpicture}[scale=0.28]
      \lgrid{-4.5}{2.5}{-5}{5.5}
      \draw[very thick, forest, domain=-3.9:1.9, samples=70, smooth] plot (\x,{\x*\x+2*\x-3});
      \fdot{forest}{-1}{-4} \fdot{redacc}{-3}{0} \fdot{redacc}{1}{0} \fdot{navy}{0}{-3}
    \end{tikzpicture}
    \end{minipage}

  \item \textbf{Contrast pair.} $p(x)=(x-2)^2-1$ and $q(x)=-(x+1)^2+4$.
    \begin{minipage}[c]{0.24\linewidth}\centering
    \begin{tikzpicture}[scale=0.26]
      \lgrid{-1.5}{5.5}{-2.5}{5.5}
      \draw[very thick, forest, domain=0.2:3.8, samples=60, smooth] plot (\x,{(\x-2)*(\x-2)-1});
      \fdot{forest}{2}{-1}
    \end{tikzpicture}\par$p$
    \end{minipage}\hfill
    \begin{minipage}[c]{0.24\linewidth}\centering
    \begin{tikzpicture}[scale=0.26]
      \lgrid{-4.5}{2.5}{-2.5}{5.5}
      \draw[very thick, redacc, domain=-3.5:1.5, samples=60, smooth] plot (\x,{-(\x+1)*(\x+1)+4});
      \fdot{redacc}{-1}{4}
    \end{tikzpicture}\par$q$
    \end{minipage}\hfill
    \begin{minipage}[c]{0.48\linewidth}
    $p$ has a \ans{minimum}: value \ans{$-1$}, at $x=$ \ans{$2$}. \\[5pt]
    $q$ has a \ans{maximum}: value \ans{$4$}, at $x=$ \ans{$-1$}.
    \end{minipage}
    \\[2pt]
    (c) \textbf{Why?} What about each graph decides which kind of value it has?
    \ansline{The direction it opens: $p$ opens up, so its vertex is the lowest point --- a minimum.}
    \ansline{$q$ opens down, so its vertex is the highest point --- a maximum. The sign of $a$ decides.}

  \item \textbf{Off a table.} A function is given only by this table.
    \begin{minipage}[c]{0.44\linewidth}\centering
    \renewcommand{\arraystretch}{1.2}
    \begin{tabular}{|c|c|c|c|c|c|}
    \hline
    $x$ & $0$ & $1$ & $2$ & $3$ & $4$ \\ \hline
    $y$ & $-3$ & $3$ & $5$ & $3$ & $-3$ \\ \hline
    \end{tabular}
    \end{minipage}\hfill
    \begin{minipage}[c]{0.52\linewidth}
    First differences: \ans{$6$} \ans{$2$} \ans{$-2$} \ans{$-6$} \\[4pt]
    Second differences: \ans{$-4$} \ans{$-4$} \ans{$-4$} --- constant? \ans{yes}
    \\[4pt]
    So the function is \ans{quadratic}. By the table's symmetry the axis is $x=$ \ans{$2$} and
    the vertex is \ans{$(2,5)$}
    \end{minipage}
\end{enumerate}
\end{notesbox}

\newpage
\begin{notesbox}{Practice, continued}
\small
\begin{enumerate}[leftmargin=1.9em, itemsep=7pt, start=4]
  \item \textbf{Special case and boundary.} How many $x$-intercepts can a parabola have?
    \ans{two, one, or none}
    \begin{minipage}[c]{0.24\linewidth}\centering
    \textbf{(a)} $y=(x-2)^2$\par
    \begin{tikzpicture}[scale=0.24]
      \lgrid{-0.5}{4.5}{-1.5}{5.5}
      \draw[very thick, forest, domain=-0.1:4.1, samples=60, smooth] plot (\x,{(\x-2)*(\x-2)});
      \fdot{forest}{2}{0}
    \end{tikzpicture}
    \end{minipage}\hfill
    \begin{minipage}[c]{0.24\linewidth}\centering
    \textbf{(b)} $y=x^2+2$\par
    \begin{tikzpicture}[scale=0.24]
      \lgrid{-2.5}{2.5}{-1.5}{5.5}
      \draw[very thick, redacc, domain=-1.8:1.8, samples=60, smooth] plot (\x,{\x*\x+2});
      \fdot{redacc}{0}{2}
    \end{tikzpicture}
    \end{minipage}\hfill
    \begin{minipage}[c]{0.48\linewidth}
    (a) zeros \ans{$1$}; the vertex sits \ans{on} the $x$-axis. \\[4pt]
    (b) zeros \ans{$0$}; the lowest output is \ans{$2$}. \\[4pt]
    (c) Opens \emph{up}, vertex \emph{above} the axis: zeros \ans{none}. Why can the answer
    never be three? \ans{it turns only once}
    \end{minipage}

  \item \textbf{Model it.} A punt sends a football up from a tee $4$ feet off the ground:
        $h(t)=-16t^2+48t+4$, with $t$ in seconds and $h$ in feet. Its graph turns at $t=1.5$.
    \begin{enumerate}[label=(\alph*), leftmargin=1.9em, itemsep=3pt]
      \item Find $h(1.5)$.
        \begin{work}
          h(1.5) &= -16(1.5)^2 + 48(1.5) + 4 \\
                 &= -36 + 72 + 4 \\
                 &= 40
        \end{work}
      \item The maximum height is \ans{$40$} ft --- an \ans{output} --- reached at $t=$
            \ans{$1.5$} s --- an \ans{input}.
      \item $h(1.5)=40$ means \ans{the ball is $40$ ft up, its highest point, $1.5$ s after the punt}
      \item The ball lands near $t=3.1$ s, but the function keeps going: as $t$ grows its outputs
            \ans{fall}. Why does the model stop making sense after it lands? \ans{height cannot be negative}
    \end{enumerate}

  \item \textbf{Multiple choice (SOL-style).} A ball's height in feet after $t$ seconds is
        $h(t)=-16t^2+96t$; its graph has vertex $(3,144)$ and zeros $t=0$ and $t=6$. What is the
        ball's \emph{maximum height}?
    \begin{enumerate}[label=\Alph*., leftmargin=2.2em, itemsep=0pt, topsep=2pt]
      \item $3$ feet
      \item \ans{$144$ feet} \hfill \textcolor{keyred}{\textbf{$\leftarrow$ correct}}
      \item $6$ feet
      \item $96$ feet
    \end{enumerate}
    Say in one sentence how you ruled out one of the others.
    \ansline{The vertex $(3,144)$ gives height $144$ ft at time $t=3$ s, so the maximum is $144$.}
    \ansline{A is $3$ --- the \emph{input}, in seconds, not a height; C is the landing time, D a coefficient.}
\end{enumerate}
\end{notesbox}

\vspace{0.08in}

\boxguard[8]
\begin{spiralbox}
\small
\textbf{Coming next --- Lesson 2.1: Graphing Quadratic Functions.} Today you read a parabola that
was already drawn; next you build one \emph{from} its equation. You will meet the same curve in
three costumes --- \textbf{standard form} $y=ax^2+bx+c$, \textbf{vertex form} $y=a(x-h)^2+k$, and
\textbf{factored form} $y=a(x-r_1)(x-r_2)$ --- each handing you a different feature for free.
\end{spiralbox}

\end{document}
